\scanpage{58}


\margin{\figure{\begin{tikzpicture}\draw (0,0) circle (.045) -- (.35,-.38);\end{tikzpicture}}}\textcircled{d} $f^*$ est conservatif, i.e. si $\lambda\in\mathrm{Fl}\,C$ est tel que $i^*\lambda$ et $j^*\lambda$ soient des isom., alors $\lambda$ est un isom.

Sous ces conditions, le foncteur naturel
\diagram{\begin{tikzcd}[ampersand replacement=\&] C\arrow[r,"\theta"] \& K \end{tikzcd}}
[où $K$ est la \deleted{\unsure{cat. des}} catégorie des triples $(G,H,u)$, $G\in\mathrm{Ob}\,C'$, $H\in\mathrm{Ob}\,C''$, $u\in\operatorname{Hom}(H,\varphi(G))$] est une \uline{équivalence de catégories}.

\uline{Ex} \deleted{\unsure{Espa}}
\[
\begin{aligned}
C&=\text{faisceaux sur }X\\
C'&\unsure{\,\cdot\,|=}\text{faisceaux sur } U\subset X\quad(U\text{ ouvert})\\
C''&=\text{faisceaux sur }Y\subset X\quad(Y=X-U)
\end{aligned}
\]

\margin{\figure{\begin{tikzpicture}\draw (0,0) circle (.035) -- (-.17,-.47);\end{tikzpicture}}}équivaut (en présence des autres axiomes) \deleted{\unsure{au fait que}} \inserted{à la condition} \deleted{diagramme} d') pour tout $F\in\mathrm{Ob}\,C$, le diagramme
\diagram{\begin{tikzcd}[ampersand replacement=\&]
F\arrow[r]\arrow[d] \& i_*i^*F\arrow[d]\\
j_*j^*F\arrow[r] \& j_*j^*i_*i^*F
\end{tikzcd}}
\diagram{\begin{tikzcd}[ampersand replacement=\&]
\text{( i.e.} \& F\arrow[r]\arrow[d] \& i_*(G)\arrow[d] \& {}\\
{} \& j_*(H)\arrow[r,no head,"\unsure{-|\to}"{marking}] \& j_*(\varphi(G)) \& \text{)}
\end{tikzcd}}
est cartésien

% archive sheet number: 57
